\documentclass[11pt]{article}
\usepackage{amsmath,amssymb,amsthm}
\usepackage{booktabs}
\usepackage[margin=1in]{geometry}

\newtheorem{theorem}{Theorem}
\newtheorem{lemma}[theorem]{Lemma}

\title{Problem 732: Standing on the Shoulders of Trolls}
\author{Project Euler Solutions}
\date{}

\begin{document}
\maketitle

\section{Problem Statement}

$N$ trolls are in a hole of depth $D_N = \frac{1}{\sqrt{2}}\sum_{n=0}^{N-1} h_n$. Each troll~$n$ has shoulder height~$h_n$, arm length~$l_n$, and IQ~$q_n$ generated by $r_n = [(5^n \bmod (10^9+7)) \bmod 101] + 50$, with $h_n = r_{3n}$, $l_n = r_{3n+1}$, $q_n = r_{3n+2}$. Trolls stack and escape from the top. Maximize total escaped IQ $Q(N)$.

Given: $Q(5) = 401$, $Q(15) = 941$. Find $Q(1000)$.

\section{Mathematical Foundation}

\begin{lemma}[Escape Condition]
A troll~$i$ atop a pile of total shoulder height~$H$ escapes if and only if $H + l_i \ge D_N$. After escape, the pile height becomes $H - h_i$.
\end{lemma}

\begin{proof}
Troll~$i$'s feet are at height $H - h_i$, shoulders at~$H$, hands at $H + l_i$. Escape requires $H + l_i \ge D_N$. Removing troll~$i$ reduces the pile by~$h_i$.
\end{proof}

\begin{theorem}[Optimal Escape Ordering]
Among all escapable trolls, the optimal order is by non-decreasing shoulder height $h_i$.
\end{theorem}

\begin{proof}
Exchange argument: suppose trolls $i, j$ are both escapable with $h_i < h_j$ and $j$ escapes before $i$. After $j$ escapes, the pile drops by~$h_j$. Had $i$ escaped first, the pile drops by $h_i < h_j$, leaving strictly more height for subsequent escapes. Hence no swap from shortest-first order improves the solution.
\end{proof}

\begin{theorem}[DP Formulation]
Sort trolls by non-decreasing $h_i$. Process in reverse order (largest~$h$ first). For each troll, either assign it to the pile (contributing $h_i$) or let it escape (contributing $q_i$ if $H + l_i \ge D_N$). The DP state is the current pile height~$H$, and the answer is $\max_H \mathrm{dp}[H]$.
\end{theorem}

\begin{proof}
By the preceding theorem the escape order is determined by the escape set. Processing from largest to smallest~$h$ and deciding pile membership yields a standard knapsack DP with state~$H$. Feasibility of each escape is checked via Lemma~1.
\end{proof}

\section{Algorithm}

\begin{enumerate}
\item Generate parameters; compute $D_N$.
\item Sort trolls by $h_i$ (non-decreasing).
\item DP in reverse sorted order: state = pile height $H$; transition = pile troll~$i$ ($H \gets H + h_i$) or escape ($+q_i$ if feasible).
\item Return $\max_H \mathrm{dp}[H]$.
\end{enumerate}

\section{Complexity Analysis}
\begin{itemize}
\item \textbf{Time:} $O(N \cdot H_{\max})$ where $H_{\max} = \sum h_i \approx 100N$; for $N=1000$ this is $\sim 10^8$.
\item \textbf{Space:} $O(H_{\max}) = O(100N)$ for the DP array.
\end{itemize}

\section{Answer}

\[
\boxed{45609}
\]

\end{document}
